% ==============================================================
% CAPÍTULO 7 — DIAGNÓSTICO E PROPOSTA DE MELHORIA
% ==============================================================

\chapter{Diagnóstico e Proposta de Melhoria}
\label{cap:diagnostico-proposta}

\section{Pontos Negativos do Processo Atual}
\label{sec:pontos-negativos}

A aplicação dos frameworks do Capítulo~\ref{cap:processo-decisorio} ao processo decisório da Capim permite consolidar cinco pontos negativos. Os três primeiros têm maior aderência direta ao relato da entrevista; os dois últimos são inferidos a partir da ausência de menção a determinados mecanismos em uma única entrevista, e devem ser tratados como hipóteses a confirmar, não como diagnósticos fechados.

\begin{enumerate}
  \item \textbf{Ausência de uma função formal de Descoberta.} Como mostrado na Seção~\ref{sec:ambidestria}, não há, no relato, evidência de hub, rito ou responsável dedicado à busca ativa de oportunidades. Os projetos mencionados (Capix, CamilA, modelo B2C) são descritos como emergindo de percepções da liderança, não de um processo sistemático de geração e captação de ideias. Pela leitura da cadeia de valor da inovação (Seção~\ref{sec:cadeia-valor}), este é o elo mais fraco do sistema.
  \item \textbf{Estrutura de avaliação igual para projetos de incerteza distinta.} O \textit{formato} do critério -- uma métrica fixada no início e formalmente revisada ao fim do trimestre -- é aplicado da mesma forma a um produto incremental maduro (sistema de gestão) e a uma aposta radical em incubação (CamilA), ainda que o \textit{valor} da métrica varie por projeto. O acompanhamento informal do andamento de cada iniciativa também ocorre na Reunião Estratégica semanal (Tabela~\ref{tab:ritos}), mas a decisão estrutural de manter, pausar ou descontinuar um projeto permanece atrelada ao ciclo trimestral, independentemente do tipo de incerteza envolvida. A tipologia de \textcite{salerno2015} sugere que projetos de categorias de processo distintas deveriam ser avaliados com estruturas de critério também distintas, não apenas com metas diferentes dentro da mesma estrutura.
  % TODO(equipe): confirmar com o Axel se a leitura de informalidade no resourcing abaixo procede.
  \item \textbf{Possível dependência de pessoas-chave e informalidade no \textit{resourcing}.} A entrevista não relata dificuldades de orçamento ou recursos para inovação, mas também não descreve mecanismos formais e documentados de alocação de tempo, orçamento ou carreira para quem lidera iniciativas radicais -- a coordenação parece concentrada na disponibilidade do PM de experimentos e da Diretoria de Produto. \textcite{salerno2025} discutem por que mecanismos formais (e não apenas informais) tendem a ser necessários para sustentar o \textit{newstream} a partir do \textit{mainstream} de forma previsível; resta confirmar com a empresa se essa informalidade é, de fato, percebida internamente como um problema.
  % TODO(equipe): Axel lembra do entrevistado mencionar uma gestão do conhecimento mais estruturada -- confirmar na transcrição/com a empresa antes de manter esta leitura.
  \item \textbf{Registro de aprendizado pouco integrado à geração de novas ideias.} A entrevista indica que projetos descontinuados são documentados (``deixa documentado lá e pronto''), o que contradiz uma leitura de ausência total de registro. O ponto mais defensável é outro: não há evidência, no relato, de que esse registro alimente de volta um funil estruturado de novas ideias -- ele parece arquivado, não reaproveitado.
  \item \textbf{Portfólio sem visão consolidada de risco e retorno.} As decisões de continuidade são tomadas projeto a projeto, sem uma visão de carteira que avalie explicitamente o balanceamento entre iniciativas de baixo e alto risco (Seção~\ref{sec:portfolio}). Na prática, isso significa que a Reunião Estratégica semanal decide sobre cada iniciativa isoladamente -- manter, pausar ou descontinuar --, sem uma pauta periódica que examine a carteira como um todo e pergunte, por exemplo, se a empresa está excessivamente concentrada em apostas de alto risco (\textit{Oysters}) e melhorias incrementais (\textit{Bread and Butter}), sem cultivar deliberadamente oportunidades de baixo risco e alto retorno (\textit{Pearls}). Esse ponto é particularmente relevante porque uma visão de portfólio explícita tende a preceder -- e a justificar -- a própria função de Descoberta proposta na Seção~\ref{sec:sistema-proposto}: sem uma leitura periódica de onde a carteira está desbalanceada, é difícil priorizar em que tipo de oportunidade essa função deveria investir esforço de busca. A ausência de iniciativas no quadrante de \textit{Pearls} é consistente com essa lacuna, mas a amostra de apenas quatro iniciativas conhecidas limita a força dessa conclusão.
\end{enumerate}

\section{Sistema Proposto}
\label{sec:sistema-proposto}

A proposta de aprimoramento preserva os elementos do processo atual que se mostraram coerentes com a literatura -- a separação física e organizacional da Garagem~2, a progressão controlada por exposição a clínicas, e a disciplina de métricas por iniciativa -- e introduz quatro mecanismos adicionais, voltados a corrigir os pontos negativos identificados.

\subsection{Função de Descoberta de Oportunidades}

Recomenda-se a criação de uma função de Descoberta, antecedendo formalmente a Garagem~2, nos termos discutidos por \textcite{salernogomes2018} e \textcite{oconnormeyer2026}. Cabe uma ressalva de proporcionalidade antes de detalhar a proposta: para uma organização de \textasciitilde150 pessoas como a Capim, criar de imediato um ``hub'' -- uma estrutura dedicada, com equipe e orçamento próprios -- pode ser desproporcional ao estágio da empresa e correr o risco de burocratizar exatamente o processo que hoje é ágil por ser informal. A alternativa mais conservadora, adotada aqui, é começar pela atribuição de um \textbf{papel} de Descoberta a uma pessoa já existente na Diretoria de Produto (não uma nova área), com escopo e tempo protegido explícitos, evoluindo para uma estrutura mais formal apenas se o volume de oportunidades identificadas justificar a dedicação de mais recursos. Concretamente, isso significa:

\begin{itemize}
  \item Designar, dentro da Diretoria de Produto, um papel explícito de ``caçador de oportunidades'' -- não necessariamente uma nova contratação ou área --, com tempo protegido para escuta ativa de clínicas, dentistas e parceiros de tecnologia (IA, fintechs, distribuidores de equipamentos odontológicos), em vez de depender apenas de percepções emergentes do PM de experimentos.
  \item Instituir um rito leve e periódico (por exemplo, mensal) de polinização cruzada entre squads, no qual cada equipe apresenta sinais de oportunidade observados em sua frente -- mecanismo direto de geração de ideias internas, na acepção de \textcite{hansenbirkinshaw2007}.
  \item Tratar cada oportunidade identificada como um item de um portfólio de descoberta, e não como uma proposta isolada a ser aprovada ou rejeitada em uma única reunião -- alinhado à recomendação de \textcite{oconnormeyer2026} de tratar domínios de inovação estratégica como portfólios de experimentos, e não como pipeline de \textit{go/kill}.
\end{itemize}

\subsection{Learning Plan para Iniciativas em Incubação}

Para projetos que ingressam na Garagem~2 com alta incerteza de mercado e/ou tecnologia -- como foi o caso da CamilA --, recomenda-se a adoção de um \textit{learning plan} formal antes da entrada nos ciclos trimestrais de avaliação tradicionais, conforme proposto por \textcite{riceoconnorpierantozzi2008} e detalhado, com casos aplicados, por \textcite{salernogomes2018}.

O \textit{learning plan} consiste em um ciclo iterativo que parte da identificação explícita do que é conhecido e do que é desconhecido em quatro categorias de incerteza -- técnica, mercadológica, organizacional e de recursos --, avalia a criticidade relativa de cada incerteza, formula hipóteses testáveis de menor custo possível e realimenta o ciclo com o aprendizado obtido, antes de qualquer compromisso financeiro maior. Esse mecanismo ocupa, no espectro de abordagens de gestão de incerteza, uma posição mais adequada a incertezas elevadas do que os \textit{stage-gates} tradicionais ou os \textit{milestones} simples -- e é particularmente recomendável nos casos em que, como na CamilA, a causa do não atingimento da métrica trimestral é a imaturidade tecnológica, e não a falta de mérito da ideia.

Na prática, recomenda-se que a transição da Garagem~2 para a Garagem~1 deixe de ser condicionada apenas ao atingimento de uma métrica de adoção fixa e passe a exigir, adicionalmente, a conclusão satisfatória de um ciclo de \textit{learning plan} -- reduzindo o risco de descontinuar prematuramente iniciativas tecnologicamente promissoras, mas ainda imaturas.

\subsection{Diferenciação Formal de Processos por Tipo de Incerteza}

Recomenda-se que a Diretoria de Produto classifique explicitamente cada iniciativa, no momento de sua entrada na Garagem~2, segundo a tipologia de \textcite{salerno2015} (Tabela~\ref{tab:tipologia}), e associe a cada tipo um critério de avaliação coerente:

\begin{itemize}
  \item Iniciativas do tipo \textit{tradicional} (mercado e tecnologia maduros, como melhorias do sistema de gestão e do Capix) continuam no regime atual de métrica trimestral e \textit{stage-gates}.
  \item Iniciativas classificadas como \textit{pausa esperando tecnologia} ou \textit{pausa esperando mercado} (como a CamilA e o modelo B2C, respectivamente) passam a ter ciclos de avaliação vinculados a marcos de maturidade tecnológica ou de validação de mercado, e não apenas ao calendário trimestral -- evitando que o relógio do trimestre force decisões binárias sobre fenômenos que evoluem em ritmo diferente.
\end{itemize}

\subsection{Mecanismos Formais de \textit{Resourcing}}

% TODO(equipe): confirmar com o Axel se mecanismos formais desse tipo fazem sentido para uma empresa do porte da Capim (~150 pessoas), ou se uma versão mais leve é mais adequada.
Inspirando-se no caso de ambidestria estrutural documentado por \textcite{salerno2025}, recomenda-se formalizar os mecanismos pelos quais squads autônomos (\textit{mainstream}) cedem tempo, orçamento e pessoas para iniciativas em incubação (\textit{newstream}): acordos explícitos de alocação de horas de engenharia por trimestre, um orçamento anual dedicado a experimentos da Garagem~2 independente do orçamento operacional dos squads, e critérios de carreira e reconhecimento específicos para quem lidera iniciativas de alta incerteza -- de modo que o sucesso nessas funções não seja avaliado pelas mesmas métricas de entrega usadas para produtos maduros. Dado o porte ainda enxuto da empresa (\textasciitilde150 colaboradores), esses acordos podem começar de forma simples -- por exemplo, um percentual fixo e pequeno da capacidade de engenharia reservado por trimestre --, evoluindo em formalidade apenas na medida em que o número de iniciativas em incubação crescer.

\subsection{Portfólio de Inovação com Revisão Periódica}

Por fim, recomenda-se que a Reunião Estratégica semanal passe a revisar, periodicamente (por exemplo, a cada trimestre), uma visão consolidada da carteira de iniciativas na matriz de risco e retorno de \textcite{goffinmitchell2010} (Tabela~\ref{tab:portfolio}), de modo a tornar explícito o balanceamento entre projetos \textit{Bread and Butter} e \textit{Oysters}, e a identificar deliberadamente a necessidade de desenvolver iniciativas no quadrante de \textit{Pearls} -- oportunidades de baixo risco e alto retorno, tipicamente amadurecidas pela própria função de Descoberta proposta na Seção anterior.

\section{Síntese da Proposta}

A Figura~\ref{fig:antes-depois} e a Tabela~\ref{tab:sintese-proposta} resumem a transição proposta do sistema atual para o sistema aprimorado.

\begin{figure}[h]
  \caption{Transição proposta: sistema atual \textit{versus} sistema aprimorado}
  \label{fig:antes-depois}
  \centering
  \begin{tikzpicture}[
      font=\footnotesize, node distance=2mm,
      now/.style={draw, fill=gray!12, rounded corners, text width=42mm, align=left, inner sep=3pt, minimum height=11mm},
      new/.style={draw, fill=blue!10, rounded corners, text width=42mm, align=left, inner sep=3pt, minimum height=11mm},
      dim/.style={font=\scriptsize\bfseries, align=center, text width=22mm}
    ]
    \node[font=\small\bfseries] (h1) at (0,1.1) {Sistema atual};
    \node[font=\small\bfseries] (h2) at (7.5,1.1) {Sistema proposto};
    \foreach \i/\d/\a/\b in {
      0/Geração de ideias/{Informal, percepções da liderança}/{Função de Descoberta + polinização},
      1/Estrutura de critério/{Métrica única, mesmo formato p/ todos}/{Estrutura diferenciada por tipo de processo},
      2/Incerteza radical/{Ciclo trimestral direto}/{\textit{Learning plan} formal},
      3/Recursos/{Alocação informal}/{Acordos formais de orçamento e RH},
      4/Portfólio/{Decisão projeto a projeto}/{Revisão periódica da carteira}} {
      \node[now] (n\i) at (0,-\i*1.3) {\a};
      \node[new] (p\i) at (7.5,-\i*1.3) {\b};
      \node[dim] at (-3.4,-\i*1.3) {\d};
      \draw[-{Stealth[length=2mm]}, gray] (n\i.east) -- (p\i.west);
    }
  \end{tikzpicture}
  \legend{Fonte: elaborado pelos autores.}
\end{figure}

\begin{table}[t]
  \centering
  \caption{Síntese: sistema atual \textit{versus} sistema proposto}
  \label{tab:sintese-proposta}
  \begin{tabularx}{\textwidth}{lXX}
    \toprule
    \textbf{Dimensão} & \textbf{Sistema atual} & \textbf{Sistema proposto} \\
    \midrule
    Geração de ideias & Informal, dependente de percepções pontuais da liderança & Função de Descoberta com escuta ativa e polinização cruzada entre squads \\
    Estrutura de critério & Métrica trimestral única, mesmo formato para todos os tipos de projeto & Estrutura diferenciada por tipo de processo (tradicional \textit{vs.} pausa) \\
    Gestão de incerteza radical & Ciclo trimestral aplicado diretamente & \textit{Learning plan} formal antes da entrada nos gates tradicionais \\
    Alocação de recursos & Informal, dependente de disponibilidade pontual & Acordos formais de orçamento, tempo de engenharia e RH dedicado \\
    Visão de portfólio & Decisões projeto a projeto & Revisão periódica da carteira na matriz de risco e retorno \\
    \bottomrule
  \end{tabularx}
  \legend{Fonte: elaborado pelos autores.}
\end{table}

A proposta não substitui o modelo atual da Capim, que já demonstra elementos sofisticados de ambidestria estrutural e de progressão controlada de risco. O objetivo é fechar especificamente a lacuna identificada na geração de ideias e tornar explícita -- em vez de implícita -- a diferenciação de governança entre inovação incremental e radical que a literatura da disciplina indica como condição para sustentar a capacidade de inovação da empresa no longo prazo.
